\section{模型假设与符号说明}

\subsection{模型假设}
为在保证问题本质不丢失的前提下建立可计算的模型，本文做出以下假设：
\begin{itemize}
	\item \textbf{假设 1（直线航段）：}任意两个任务节点之间采用两点间的水平直线作为运输航段，水平巡航距离按局部平面投影（WGS84）计算。
	\item \textbf{假设 2（巡航海拔）：}计划巡航海拔取该航段所经过 DEM 像元的最高地面高程以上 50\,m，且不低于起终点作业高度；爬升与下降为竖直运动，水平巡航为定高飞行。
	\item \textbf{假设 3（能耗口径）：}水平巡航能耗由等效航程反推（$E^{\mathrm{hor}}=E^{\mathrm{use}}\cdot d/L_g(q)$），爬升附加能耗取重力势能除以爬升能耗效率；下降不单独计能耗（效率为 0）。
	\item \textbf{假设 4（理想环境）：}忽略风、雨、能见度等气象因素对飞行时间与能耗的影响；无人机起降、悬停、交接过程稳定可靠。
	\item \textbf{假设 5（充电模型）：}电池与能源组件统一采用题给两阶段等效充电模型，多块电池可并行充电，同一资源占用与充电时段不重叠。
	\item \textbf{假设 6（中继共享）：}一架中继无人机可同时为多架位于其覆盖范围内的运输无人机提供通信保障，但不允许多跳中继转发。
	\item \textbf{假设 7（通信判定）：}通信链路按自由空间传播损耗加地形遮挡附加损耗进行链路预算，双向取较小门限判定链路可用性。
\end{itemize}

\subsection{符号说明}
本文使用的主要符号及其含义如表~\ref{tab:symbols}所示。

\begin{table}[htbp]
	\caption{符号说明}
	\label{tab:symbols}
	\centering
	\begin{tabular}{cl}
		\toprule
		符号 & 说明 \\
		\midrule
		$g\in\{A,B,C\}$ & 运输无人机机型编号 \\
		$S_i$（$i=1,\dots,15$） & 第 $i$ 个服务区 \\
		$O01$ & 临时调度中心 \\
		$Q_g$ & 机型 $g$ 的最大载货质量（kg） \\
		$q$ & 当前有效载荷（kg） \\
		$L_{g0},L_{gF}$ & 机型 $g$ 的空载/满载标准航程（m） \\
		$L_g(q)$ & 机型 $g$ 携带载荷 $q$ 的等效航程（m） \\
		$E_g^{\mathrm{use}}$ & 机型 $g$ 单组电池可用能量（kWh） \\
		$\rho_g$ & 机型 $g$ 的返航安全余量比例 \\
		$d_{ij}$ & 节点 $i\to j$ 的水平巡航距离（m） \\
		$h^{+}_{ij},h^{-}_{ij}$ & 航段 $i\to j$ 的爬升/下降高度（m） \\
		$v^{\uparrow}_g,v^{c}_g,v^{\downarrow}_g$ & 机型 $g$ 的爬升/巡航/下降速度（m/s） \\
		$E_{gij}(q)$ & 机型 $g$ 航段 $i\to j$ 载荷 $q$ 的能耗（kWh） \\
		$t_{gij}$ & 机型 $g$ 航段 $i\to j$ 的飞行时间（s） \\
		$b_{ijt}$ & 地形遮挡变量（0/1） \\
		$L^{\max}_{a\leftrightarrow b}$ & 双向链路最大允许传播损耗（dB） \\
		$A_{ijt}$ & 链路可用性变量 \\
		$s$ & 电池/能源组件剩余荷电状态（SOC） \\
		\midrule
		$s_p$ & 运输架次 $p$ 的开始时刻（问题三的错峰决策变量） \\
		$\pi$ & 派工优先级排列（问题二解的第二分量） \\
		$y_r\in\{0,1\}$ & 悬停站候选点 $r$ 是否启用（选址二元变量） \\
		$h_r$ & 候选点 $r$ 的悬停离地高度（m），$h_r\in\{50,100,\dots,300\}$ \\
		$a_{jr}$ & 覆盖系数：候选点 $r$ 能否完整覆盖失效区间 $j$ \\
		$C_p(t)$ & 架次 $p$ 在 $t$ 时刻的通信状态（1 为连通） \\
		$w_i$ & 时间积分权重 $w_i=(t_{i+1}-t_{i-1})/2$ \\
		$K$ & 架次数上限（问题二）或任务组数（问题四） \\
		$M$ & big-$M$ 析取约束中的足够大常数，须满足 $M\ge UB+\max_c\tau_c$ \\
		$\kappa$ & 电池可用能量折减系数，$E_g^{\mathrm{use}}\to\kappa E_g^{\mathrm{use}}$ \\
		\bottomrule
	\end{tabular}
\end{table}
